\documentclass[english, parskip=full]{scrartcl}
\usepackage[english]{babel}  % Sprache auf Deutsch 
\usepackage[utf8]{inputenc}  % Kodierung der Datei
\usepackage[left=2cm, right=2cm, top=2cm, bottom=3cm, footskip=1.5cm]{geometry}
\usepackage{color}
\usepackage{graphicx, gensymb}
\usepackage{amsfonts, cancel, amssymb}
\usepackage{amsmath, amsthm, array, esint}
\usepackage{bm} % bold math symbols, e.g. \bm{\sigma} etc.
\usepackage{booktabs, multirow}  % schönere Tabellen
%\usepackage{scrlayer-scrpage}    % Manipulation des Seitenstils
%\pagestyle{scrheadings}  % Stil für die Seite setzen
\usepackage{tabularx}
\usepackage{setspace}
%\automark{section}  % Abschnittsnamen als Seitenbeschriftung 
%\setheadsepline{.5pt}    % Kopzeile durch Linie abtrennen
\usepackage[hidelinks]{hyperref}  % Links und weitere PDF-Features
\usepackage{typearea, tikz, pgffor, verbatim}
\usetikzlibrary{calc, arrows.meta,arrows, graphs, quotes, positioning, angles, babel}
\usepackage[cache=false, outputdir=build]{minted}
\usemintedstyle{vs}
\usepackage{placeins} % provides \FloatBarrier

\definecolor{bg}{rgb}{0.9,0.9,0.9}
\newcommand\numberthis{\addtocounter{equation}{1}\tag{\theequation}}
\usepackage{svg}
\usepackage{siunitx}
\usepackage{physics}
\usepackage{tensor}
\usepackage{slashed}
\usepackage{bbm}
\usepackage[compat=1.1.0]{tikz-feynman}

\usepackage{caption}
\captionsetup{
    % width=.8\textwidth,
    % indention=1cm,
    margin=0.5cm,
    format=plain,
    labelfont=bf
}
%\usepackage{subcaption}
%\subcaptionsetup{
%    indention=0cm,
%    format=hang,
%    margin=0.5cm,
%    labelfont=normalfont
%}
\usepackage[english,capitalise]{cleveref}
\usepackage[
    backend=biber,
    sorting=none,
    style=phys,
    giveninits=true,
    sortcites=true,
    citestyle=numeric-comp,
    % url=false,
    % doi=false,
    biblabel=brackets,
    % eprint=false,
    maxcitenames=3]{biblatex}
\usepackage{csquotes}
%\addbibresource{bibliography.bib}

\usepackage{mathtools}
% use \abs for absolute values
% use \abs for \left ... \right version and \abs[\bigg for \bigg version etc.
\let\abs\undefined
\let\bra\undefined
\let\braket\undefined
\DeclarePairedDelimiter\abs{\vert}{\vert}
\DeclarePairedDelimiter\kla{(}{)}
\DeclarePairedDelimiter\klb{[}{]}
\DeclarePairedDelimiter\klc{\{}{\}}
\DeclarePairedDelimiter\bra{\langle}{\rangle}
\DeclarePairedDelimiterX\brb[2]{\langle #1}{\rangle}{\delimsize\vert #2 }
\DeclarePairedDelimiterX\brc[3]{\langle #1}{\rangle}{\delimsize\vert #2 \delimsize\vert #3 }

\renewcommand{\i}{\text{i}}
\newcommand{\e}{\text{e}}
\DeclareMathOperator{\sgn}{sgn}
\newcommand{\kom}[1]{\overset{\substack{#1 \\ |}}{=}}
\newcommand{\koml}[2]{\overset{\substack{#1 \\ |}}{#2}}
\newcommand{\intx}[4]{\int_{#1}^{#2} #3 \,\mathrm{d}#4}
\newcommand{\intr}[2]{\int_{\mathbb{R}} #1 \, \mathrm{d}#2}
\newcommand{\intrnd}[3]{\int_{\mathbb{R}^{#1}} #2 \, \mathrm{d}^{#1}#3}
\newcommand{\intrpi}[2]{\int_{\mathbb{R}} #1 \, \frac{\mathrm{d}#2}{2\pi}}
\newcommand{\intrndpi}[3]{\int_{\mathbb{R}^{#1}} #2 \, \frac{\mathrm{d}^{#1}#3}{(2\pi)^{#1}}}
\newcommand{\oper}[2]{\vert #1 \rangle \langle #2 \vert}
\newcommand{\Text}[1]{\text{#1}}    % this is relevant because \text does not autocomplete to the default '\text{@}' but rather '\textbf' etc.
\newcommand{\powe}[1]{\e^{#1}}
\newcommand{\pow}[2]{{#1}^{#2}}
\newcommand{\Ket}[1]{\vert #1 \rangle}
\newcommand{\Bra}[1]{\langle #1 \vert}
\newcommand*{\Fourier}{\textsc{Fourier}\ }
\newcommand*{\F}{\mathcal{F}}
\DeclareMathOperator{\arsinh}{arsinh}
\DeclareMathOperator{\arcosh}{arcosh}
\DeclareMathOperator{\artanh}{artanh}
\DeclareMathOperator{\sinc}{sinc}
\DeclareMathOperator{\cscc}{cscc}
\DeclareMathOperator{\sinhc}{sinhc}
\DeclareMathOperator{\cschc}{cschc}
\newcommand{\conv}{\mathop{\scalebox{3.5}{\raisebox{-0.38ex}{\makebox[0.5\width][c]{$\ast$}}}}\limits}
\newcommand*{\unity}{\text{\usefont{U}{bbold}{m}{n}1}}   % operator one from :: https://tex.stackexchange.com/questions/566780/import-mathbb1-character-from-the-unicode-math-package

\renewcommand{\d}{\text{d}}
\renewcommand{\r}{\text{r}}
\renewcommand{\t}{\text{t}}
\renewcommand{\a}{\text{a}}
\renewcommand{\o}{\text{o}}
\renewcommand{\F}{\text{F}}
\renewcommand{\c}{\text{c}}
\newcommand{\A}{\text{A}}
\newcommand{\B}{\text{B}}
\newcommand{\R}{\text{R}}
\newcommand{\M}{\text{M}}
\newcommand{\m}{\text{m}}
\newcommand{\s}{\text{s}}
\newcommand{\f}{\text{f}}
\newcommand{\K}{\text{K}}
\newcommand{\hc}{\text{h.c.}}
\newcommand{\kf}{k_\F}
\newcommand{\vf}{v_\F}
\newcommand{\const}{\text{const.}}
\renewcommand{\L}{\text{L}}
\newcommand{\gray}[1]{\bgroup\color{white!40!black}#1\egroup}
\newcommand{\TODO}[1]{\bgroup\color{red!60!black}#1\egroup}
\newcommand\QnA[1]{\bgroup\color{blue}#1\egroup}
\newcommand\highlight[1]{\bgroup\color{green!60!black}#1\egroup}

\newcommand{\cabx}[3]{\begin{smallmatrix}\gray{A}&\gray{:}&#1 \\ \gray{B}&\gray{:}&#2 \\ \gray{\Xi}&\gray{:}&#3\end{smallmatrix}}
\newcommand{\zabx}[3]{\begin{smallmatrix}\gray{A}&\gray{:}&#1 \\ \gray{B}&\gray{:}&#2 \\ \gray{Z}&\gray{:}&#3\end{smallmatrix}}
\newcommand{\abx}[3]{\begin{smallmatrix}#1 \\ #2 \\ #3\end{smallmatrix}}

\newcommand{\normord}[1]{
  {:\mathrel{\mspace{1mu}#1\mspace{1mu}}:}
}
\newcommand{\varnormord}[1]{
  {\mathrel{\mspace{1mu}\substack{\ast\\[-1.5pt] \ast}#1\substack{\ast\\[-1.5pt] \ast}\mspace{1mu}}}
}

% punctuation styles (see https://tex.stackexchange.com/questions/56063/how-to-properly-typeset-footnotes-superscripts-after-punctuation-marks)
\let\origfootnote\footnote
\renewcommand{\footnote}[1]{\kern.06em\origfootnote{#1}}
\newcommand{\punctfootnote}[1]{\kern-.06em\origfootnote{#1}}

% Multiline equations
\newcommand{\bsplit}[1]{\begin{aligned}[t]#1\end{aligned}}

\newcommand*{\titel}{Functional Determinants in QFT}
\newcommand*{\autor}{Joachim Schwardt}

\hypersetup{pdfauthor={\autor}, pdftitle={\titel}}

%\titlehead{titlehead}
%\subject{SUBJECT}
\title{\titel}
\author{\autor}
\date{\today}

\begin{document}

%\begin{titlepage}
\maketitle  % Titel setzen
%\tableofcontents  % Inhaltsverzeichnis setzen
%\end{titlepage}


%%% THIS IS ~FunctionalDeterminants_QFT_Dunne !!!

\section{Introduction}
Consider the source-less bosonic action
\begin{align}
S[\phi] &= \intx{ }{ }{\phi \kla*{-\square + V}\phi}{x}
\end{align}
with a scalar field $\phi$.
Then the generating functional is defined as
\begin{align}
\mathcal{Z} &= \int \powe{-S[\phi]} \mathcal{D}\phi \equiv \mathcal{N} \sqrt{\frac{\pi}{\det \kla*{-\square + V}}}
\end{align}
where we defined the functional determinant by analogy to gauss integrals, up to some normalization constant $\mathcal{N}$.
Expanding the action about the classical solution gives
\begin{align}
S[\phi] &= S[\phi_\text{cl}] + \intx{ }{ }{\phi(x) \frac{\delta S}{\delta \phi(x)} \Big\vert_{\phi = \phi_\text{cl}}}{x} + \frac{1}{2} \intx{ }{ }{\intx{ }{ }{\phi(x) \frac{\delta^2 S}{\delta\phi(x) \delta\phi(y)} \Big\vert_{\phi = \phi_\text{cl}}}{y}}{x} + \dots
\end{align}
and defining the second order derivative as $K(x,y)$ yields
\begin{align}
\mathcal{Z} &\approx \mathcal{N} \powe{-S[\phi_\text{cl}]} \sqrt{\frac{\pi}{\det K}}.
\end{align}
\section{$\zeta$-function Regularization}
Consider an eigenvalue equation for some differential operator
\begin{align}
M\phi_n &= \lambda_n\phi_n.
\end{align}
Define the corresponding $\zeta$-function by
\begin{align}
\zeta(s) &= \tr\klc*{\frac{1}{M^s}} = \sum_n \frac{1}{\lambda_n^s}.
\end{align}
Then we find the relations
\begin{align}
\zeta'(s) &= -\sum_n \frac{\log(\lambda_n)}{\lambda_n^s},\\
\zeta'(0) &= -\sum_n\log(\lambda_n) = -\log\kla*{\prod_n \lambda_n}
\end{align}
which allows us to define the functional determinant as $\det M = \powe{-\zeta'(0)}$.
\section{Borel summation}
Consider the alternating divergent series
\begin{align}
f(g) &= \sum_{n\ge 0} (-1)^n n! g^n.
\end{align}
Define a new series by representing $n!$ via the gamma function and interchange integration and summation
\begin{align}
\tilde{f}(g) &= \intx{0}{\infty}{ \sum_{n\ge 0} \powe{-t} t^n (-1)^n g^n }{t} = \frac{1}{g} \intx{0}{\infty}{\frac{\powe{-\frac{x}{g}}}{1 + x}}{x} \equiv \textsc{Borel}\text{-sum of }f(g).
\end{align}
More generally, we can compute the asymptotic series expansion
\begin{align}
f(g) \sim \sum_{n\ge 0} (-1)^n \Gamma\kla*{b n + c} g^n \equiv \frac{1}{b} \intx{0}{\infty}{ \frac{1}{x(1+x)} \kla*{\frac{x}{g}}^\frac{c}{b} \powe{-\kla*{\frac{x}{g}}^{1/b}} }{x}.
\end{align}
\section{The Gel'fand-Yaglom formalism}
Assume that we do not know the explicit eigenvalues of the differential operator $M$, but instead are given some function satisfying $F(\lambda_n) = 0$ for all $n$.
Then we can write
\begin{align}
\zeta(s) &= \frac{1}{2\pi\i} \oint_C \lambda^{-s} \frac{\d}{\d\lambda} \log F(\lambda) \,\d\lambda\\
&= \frac{\sin(\pi s)}{\pi} \intx{0}{-\infty}{\lambda^{-s}\frac{\d}{\d\lambda} \log F(\lambda) }{\lambda}.
\end{align}
Then using the definition of the functional determinant we find
\begin{align}
\det M &= \powe{-\zeta'(0)} = \powe{-\log F(-\infty) + \log F(0)} = \frac{F(0)}{F(-\infty)}.
\end{align}
Usually one may assume $F(-\infty)$ to be independent of the details of the potential $V$ and thus
\begin{align}
\frac{\det M}{\det M_\text{free}} &= \frac{F(0)}{F_\text{free}(0)}.
\end{align}
\subsection{One-dimensional Schrödinger operators}
Consider the operator $M=-\frac{\d^2}{\d x^2} + V(x)$ on the interval $x\in[0,1]$.
Suppose that we want to solve the eigenvalue problem $M\phi_n = \lambda_n\phi_n$ with $\phi_n(0)=0=\phi_n(1)$ where $0<\lambda_1<\lambda_2<\dots$.
Now let
\begin{align}
Mu_\lambda &= \lambda u_\lambda
\end{align}
with $u_\lambda(0) = 0$ and $u_\lambda'(0) = 1$.
By definition we have $u_{\lambda_n}(1) = 0$ so we set $F(\lambda) = u_\lambda(1)$.
Then we find
\begin{align}
\frac{\det M}{\det M_\text{free}} &= \frac{u_0(1)}{u_0^\text{free}(1)}.
\end{align}
The free problem is simple because
\begin{align}
{u_0^\text{free}}''(x) &= 0
\end{align}
has the immediate solution $u_0^\text{free}(x) = x$.
\subsection{$1+1$-dimensional QFT}
Consider the scalar Lagrangian
\begin{align}
\mathcal{L} &= \frac{1}{2} \partial_\mu\phi\partial^\mu\phi - U(\phi)
\end{align}
with the potential $U(\phi) =\frac{1}{g^2} \kla*{1 - \cos(g\phi)}$.
We want to find a semi-classical solution and hence first consider the energy
\begin{align}
E &= \intx{ }{ }{\kla*{\frac{1}{2}\phi'^2 + U(\phi)}}{x}
\end{align}
that is minimised at $\phi' = \sqrt{2U}$.
This differential equation can be integrated via
\begin{align}
ax+b &= \intx{ }{ }{\frac{\phi'(x)}{\sqrt{2U}}}{x} = \intx{ }{ }{\frac{g}{\sqrt{2 - 2\cos(g\phi)}}}{\phi} = \frac{g}{2} \intx{ }{ }{\csc\kla*{\frac{g}{2}\phi}}{\phi} = \log\tan \kla*{\frac{g}{4}\phi}.
\end{align}
Solving for the field yields $\phi(x) = \frac{4}{g} \arctan\kla*{\powe{a(x-x_0)}}$.


%\begin{thebibliography}{9}
%\bibitem{example} description, \url{https://www.exmapleurl}, day.\,month~2021.
%\end{thebibliography}

\end{document}